%% chapters/00-kata-pengantar.tex
%% Kata Pengantar
%% Aturan: .cursor/rules/disertation/12-abstrak-front-matter.mdc

\begin{katapengantar}

Puji dan syukur penulis panjatkan kepada Tuhan Yang Maha Esa atas
rahmat dan karunia-Nya sehingga disertasi ini dapat diselesaikan.
Disertasi ini merupakan salah satu syarat untuk memperoleh gelar
Doktor di Program Studi Doktor Teknik dan Manajemen Industri, Fakultas
Teknologi Industri, Institut Teknologi Bandung.

Penyelesaian disertasi ini tidak lepas dari dukungan dan bimbingan
berbagai pihak.
Penulis mengucapkan terima kasih yang sebesar-besarnya kepada:

\begin{enumerate}
  \item \textbf{Prof. Dr. Ir. Kadarsah Suryadi, DEA}, selaku Promotor,
        atas bimbingan, arahan ilmiah, dan dukungan moril yang tak
        henti-hentinya selama proses penelitian dan penulisan disertasi ini.
  \item \textbf{Prof. Ir. Bermawi Priyatna Iskandar, M.Sc., Ph.D.},
        selaku Anggota Tim Pembimbing 1, atas masukan berharga dalam
        pengembangan metodologi dan kerangka analisis penelitian.
  \item \textbf{Prof. Dr. Ir. Bambang Riyanto Trilaksono}, selaku
        Anggota Tim Pembimbing 2, atas wawasan teknis dan diskusi
        mendalam dalam bidang sistem kendali dan kecerdasan buatan
        yang memperkaya penelitian ini.
  \item Pimpinan dan segenap civitas akademika Institut Teknologi Bandung,
        atas fasilitas dan dukungan akademik yang disediakan.
  \item Rekan-rekan seperjuangan di Program Studi Doktor Teknik dan
        Manajemen Industri, atas diskusi dan semangat yang saling
        menguatkan.
  \item Keluarga tercinta, atas doa, dukungan, dan pengertian tanpa
        batas selama perjalanan studi doktoral ini.
\end{enumerate}

Penulis menyadari bahwa disertasi ini masih jauh dari sempurna.
Saran dan kritik yang konstruktif sangat diharapkan demi perbaikan
dan pengembangan lebih lanjut.
Semoga disertasi ini memberikan manfaat bagi perkembangan ilmu
pengetahuan dan teknologi, khususnya di bidang perawatan prediktif
berbasis kecerdasan buatan dan pembelajaran mesin untuk sistem produksi
industri di Indonesia.

\vspace{1.5cm}
\begin{flushright}
Bandung, Mei 2026

\vspace{1.5cm}
Toto Suharto
\end{flushright}

\end{katapengantar}
